\documentclass[11pt]{article}
\usepackage[margin=2.2cm]{geometry}
\usepackage{amsmath,amssymb,amsthm,booktabs,enumitem}
\usepackage[hidelinks]{hyperref}
\title{Peer review: Rick's \emph{Two indecomposability theorems for 0-Hecke modules:
induced simples (Q72(i)) and Kim--Lee--Oh Conjecture 7.5 (Q73)}}
\author{Clio}
\date{2026-11-12}
\makeatletter
%% MIRROR of proofs/protocol-header.tex (clio-vega/proofs), copied 2026-10-16.
%% reviews/ is a SEPARATE REPO (clio-vega/rick-review) and must build standalone for
%% Rick and Robin, so a symlink or TEXINPUTS dependency would break for them. Hence a
%% copy -- but a VERIFIABLE one, because a silent fork is the same fault in new clothes:
%%   source md5 at copy time: f828569e31971b23bc5726db5d23ecef
%%   drift check:  md5sum ../proofs/protocol-header.tex   # compare with the line above
%% If they differ, re-copy from proofs/ -- that file is the source of truth.
%%
%% protocol-header.tex -- PROTOCOL §2.3 page-1 block, as a TEMPLATE not an instance.
%%
%% PROTOCOL §2.3: "Every PDF must state, on its first page: author, date,
%% recipient, a title, and the commit hash of the work-in-progress push it
%% corresponds to.  A PDF with no commit hash cannot be checked against the
%% record and should be treated as a draft opinion, not a result."
%%
%% Written 2026-10-16 because the fault had recurred twice and I had twice
%% fixed only the instance:
%%   - 2026-10-14: the Q473 PDF went to Rick with no recipient and no hash.
%%   - 2026-10-14: the Q460 PDF went to Rick saying "Recipient: Robin Langer".
%%     He noticed, and took the commit hash as authority instead.
%% Both times I repaired the one file and wrote "fix the template, not just the
%% instance" in my notes.  This file is that fix.  \input it; do not retype it.
%%
%% USAGE -- TWO PLACES, NOT ONE.  The \input goes in the PREAMBLE inside
%% \makeatletter; only \protocolheader goes after \maketitle:
%%
%%   \makeatletter              <- REQUIRED: \protocol@date has an @ in its name
%%   %% MIRROR of proofs/protocol-header.tex (clio-vega/proofs), copied 2026-10-16.
%% reviews/ is a SEPARATE REPO (clio-vega/rick-review) and must build standalone for
%% Rick and Robin, so a symlink or TEXINPUTS dependency would break for them. Hence a
%% copy -- but a VERIFIABLE one, because a silent fork is the same fault in new clothes:
%%   source md5 at copy time: f828569e31971b23bc5726db5d23ecef
%%   drift check:  md5sum ../proofs/protocol-header.tex   # compare with the line above
%% If they differ, re-copy from proofs/ -- that file is the source of truth.
%%
%% protocol-header.tex -- PROTOCOL §2.3 page-1 block, as a TEMPLATE not an instance.
%%
%% PROTOCOL §2.3: "Every PDF must state, on its first page: author, date,
%% recipient, a title, and the commit hash of the work-in-progress push it
%% corresponds to.  A PDF with no commit hash cannot be checked against the
%% record and should be treated as a draft opinion, not a result."
%%
%% Written 2026-10-16 because the fault had recurred twice and I had twice
%% fixed only the instance:
%%   - 2026-10-14: the Q473 PDF went to Rick with no recipient and no hash.
%%   - 2026-10-14: the Q460 PDF went to Rick saying "Recipient: Robin Langer".
%%     He noticed, and took the commit hash as authority instead.
%% Both times I repaired the one file and wrote "fix the template, not just the
%% instance" in my notes.  This file is that fix.  \input it; do not retype it.
%%
%% USAGE -- TWO PLACES, NOT ONE.  The \input goes in the PREAMBLE inside
%% \makeatletter; only \protocolheader goes after \maketitle:
%%
%%   \makeatletter              <- REQUIRED: \protocol@date has an @ in its name
%%   %% MIRROR of proofs/protocol-header.tex (clio-vega/proofs), copied 2026-10-16.
%% reviews/ is a SEPARATE REPO (clio-vega/rick-review) and must build standalone for
%% Rick and Robin, so a symlink or TEXINPUTS dependency would break for them. Hence a
%% copy -- but a VERIFIABLE one, because a silent fork is the same fault in new clothes:
%%   source md5 at copy time: f828569e31971b23bc5726db5d23ecef
%%   drift check:  md5sum ../proofs/protocol-header.tex   # compare with the line above
%% If they differ, re-copy from proofs/ -- that file is the source of truth.
%%
%% protocol-header.tex -- PROTOCOL §2.3 page-1 block, as a TEMPLATE not an instance.
%%
%% PROTOCOL §2.3: "Every PDF must state, on its first page: author, date,
%% recipient, a title, and the commit hash of the work-in-progress push it
%% corresponds to.  A PDF with no commit hash cannot be checked against the
%% record and should be treated as a draft opinion, not a result."
%%
%% Written 2026-10-16 because the fault had recurred twice and I had twice
%% fixed only the instance:
%%   - 2026-10-14: the Q473 PDF went to Rick with no recipient and no hash.
%%   - 2026-10-14: the Q460 PDF went to Rick saying "Recipient: Robin Langer".
%%     He noticed, and took the commit hash as authority instead.
%% Both times I repaired the one file and wrote "fix the template, not just the
%% instance" in my notes.  This file is that fix.  \input it; do not retype it.
%%
%% USAGE -- TWO PLACES, NOT ONE.  The \input goes in the PREAMBLE inside
%% \makeatletter; only \protocolheader goes after \maketitle:
%%
%%   \makeatletter              <- REQUIRED: \protocol@date has an @ in its name
%%   \input{protocol-header}    <- PREAMBLE, while \@date is still live
%%   \makeatother
%%   ...
%%   \begin{document}
%%   \maketitle
%%   \protocolheader{Rick (grandpa-rick)}{clio-vega/proofs}{60e4ce3}
%%
%% CORRECTED 2026-10-16 (peer review): this block previously read "USAGE,
%% immediately after \maketitle: \input{protocol-header} \protocolheader{...}",
%% which is WRONG TWICE -- it omits \makeatletter, and it puts the \input after
%% \maketitle, by which point \@date has been cleared.  I followed it verbatim
%% while building the PROVE-241 review and page 1 printed a literal "@@date"
%% plus a stray "datedate", with pdflatex exiting 0.  The three proofs/*.tex that
%% use this file were all written with the CORRECT pattern (\makeatletter +
%% \input in the preamble), so their PDFs are fine and were verified by
%% pdftotext -- the defect was in this instruction, not in the macro.
%% An instruction that encodes a fault renews it in every session that obeys it.
%%
%% The recipient argument is MANDATORY and has no default on purpose: a default
%% is what produced "Recipient: Robin Langer" on a PDF addressed to Rick.

%% \maketitle CLEARS \@date (and \@author, \@title) in the article class, so a
%% macro called AFTER \maketitle that reads \@date prints an EMPTY row.  Caught
%% 2026-10-16 only because I ran pdftotext on page 1 instead of trusting rc=0:
%% the build was clean and the Date field was blank.  Hence: snapshot the date
%% HERE, at \input time in the preamble, where \@date is still live.
\let\protocol@date\@date

\newcommand{\protocolheader}[3]{%
  \begin{center}
  \fbox{\begin{minipage}{0.92\textwidth}
  \small
  \textbf{PROTOCOL §2.3 record.}\\[1mm]
  \begin{tabular}{@{}ll@{}}
  Author     & Clio \\
  Date       & \protocol@date \\
  Recipient  & #1 \\
  Repository & \texttt{#2} \\
  Commit     & \texttt{#3} \\
  \end{tabular}\\[1mm]
  This PDF corresponds to the pushed commit above.  Where this text and the
  commit disagree, \textbf{the commit is authoritative}.
  \end{minipage}}
  \end{center}
  \vspace{2mm}%
}
    <- PREAMBLE, while \@date is still live
%%   \makeatother
%%   ...
%%   \begin{document}
%%   \maketitle
%%   \protocolheader{Rick (grandpa-rick)}{clio-vega/proofs}{60e4ce3}
%%
%% CORRECTED 2026-10-16 (peer review): this block previously read "USAGE,
%% immediately after \maketitle: %% MIRROR of proofs/protocol-header.tex (clio-vega/proofs), copied 2026-10-16.
%% reviews/ is a SEPARATE REPO (clio-vega/rick-review) and must build standalone for
%% Rick and Robin, so a symlink or TEXINPUTS dependency would break for them. Hence a
%% copy -- but a VERIFIABLE one, because a silent fork is the same fault in new clothes:
%%   source md5 at copy time: f828569e31971b23bc5726db5d23ecef
%%   drift check:  md5sum ../proofs/protocol-header.tex   # compare with the line above
%% If they differ, re-copy from proofs/ -- that file is the source of truth.
%%
%% protocol-header.tex -- PROTOCOL §2.3 page-1 block, as a TEMPLATE not an instance.
%%
%% PROTOCOL §2.3: "Every PDF must state, on its first page: author, date,
%% recipient, a title, and the commit hash of the work-in-progress push it
%% corresponds to.  A PDF with no commit hash cannot be checked against the
%% record and should be treated as a draft opinion, not a result."
%%
%% Written 2026-10-16 because the fault had recurred twice and I had twice
%% fixed only the instance:
%%   - 2026-10-14: the Q473 PDF went to Rick with no recipient and no hash.
%%   - 2026-10-14: the Q460 PDF went to Rick saying "Recipient: Robin Langer".
%%     He noticed, and took the commit hash as authority instead.
%% Both times I repaired the one file and wrote "fix the template, not just the
%% instance" in my notes.  This file is that fix.  \input it; do not retype it.
%%
%% USAGE -- TWO PLACES, NOT ONE.  The \input goes in the PREAMBLE inside
%% \makeatletter; only \protocolheader goes after \maketitle:
%%
%%   \makeatletter              <- REQUIRED: \protocol@date has an @ in its name
%%   \input{protocol-header}    <- PREAMBLE, while \@date is still live
%%   \makeatother
%%   ...
%%   \begin{document}
%%   \maketitle
%%   \protocolheader{Rick (grandpa-rick)}{clio-vega/proofs}{60e4ce3}
%%
%% CORRECTED 2026-10-16 (peer review): this block previously read "USAGE,
%% immediately after \maketitle: \input{protocol-header} \protocolheader{...}",
%% which is WRONG TWICE -- it omits \makeatletter, and it puts the \input after
%% \maketitle, by which point \@date has been cleared.  I followed it verbatim
%% while building the PROVE-241 review and page 1 printed a literal "@@date"
%% plus a stray "datedate", with pdflatex exiting 0.  The three proofs/*.tex that
%% use this file were all written with the CORRECT pattern (\makeatletter +
%% \input in the preamble), so their PDFs are fine and were verified by
%% pdftotext -- the defect was in this instruction, not in the macro.
%% An instruction that encodes a fault renews it in every session that obeys it.
%%
%% The recipient argument is MANDATORY and has no default on purpose: a default
%% is what produced "Recipient: Robin Langer" on a PDF addressed to Rick.

%% \maketitle CLEARS \@date (and \@author, \@title) in the article class, so a
%% macro called AFTER \maketitle that reads \@date prints an EMPTY row.  Caught
%% 2026-10-16 only because I ran pdftotext on page 1 instead of trusting rc=0:
%% the build was clean and the Date field was blank.  Hence: snapshot the date
%% HERE, at \input time in the preamble, where \@date is still live.
\let\protocol@date\@date

\newcommand{\protocolheader}[3]{%
  \begin{center}
  \fbox{\begin{minipage}{0.92\textwidth}
  \small
  \textbf{PROTOCOL §2.3 record.}\\[1mm]
  \begin{tabular}{@{}ll@{}}
  Author     & Clio \\
  Date       & \protocol@date \\
  Recipient  & #1 \\
  Repository & \texttt{#2} \\
  Commit     & \texttt{#3} \\
  \end{tabular}\\[1mm]
  This PDF corresponds to the pushed commit above.  Where this text and the
  commit disagree, \textbf{the commit is authoritative}.
  \end{minipage}}
  \end{center}
  \vspace{2mm}%
}
 \protocolheader{...}",
%% which is WRONG TWICE -- it omits \makeatletter, and it puts the \input after
%% \maketitle, by which point \@date has been cleared.  I followed it verbatim
%% while building the PROVE-241 review and page 1 printed a literal "@@date"
%% plus a stray "datedate", with pdflatex exiting 0.  The three proofs/*.tex that
%% use this file were all written with the CORRECT pattern (\makeatletter +
%% \input in the preamble), so their PDFs are fine and were verified by
%% pdftotext -- the defect was in this instruction, not in the macro.
%% An instruction that encodes a fault renews it in every session that obeys it.
%%
%% The recipient argument is MANDATORY and has no default on purpose: a default
%% is what produced "Recipient: Robin Langer" on a PDF addressed to Rick.

%% \maketitle CLEARS \@date (and \@author, \@title) in the article class, so a
%% macro called AFTER \maketitle that reads \@date prints an EMPTY row.  Caught
%% 2026-10-16 only because I ran pdftotext on page 1 instead of trusting rc=0:
%% the build was clean and the Date field was blank.  Hence: snapshot the date
%% HERE, at \input time in the preamble, where \@date is still live.
\let\protocol@date\@date

\newcommand{\protocolheader}[3]{%
  \begin{center}
  \fbox{\begin{minipage}{0.92\textwidth}
  \small
  \textbf{PROTOCOL §2.3 record.}\\[1mm]
  \begin{tabular}{@{}ll@{}}
  Author     & Clio \\
  Date       & \protocol@date \\
  Recipient  & #1 \\
  Repository & \texttt{#2} \\
  Commit     & \texttt{#3} \\
  \end{tabular}\\[1mm]
  This PDF corresponds to the pushed commit above.  Where this text and the
  commit disagree, \textbf{the commit is authoritative}.
  \end{minipage}}
  \end{center}
  \vspace{2mm}%
}
    <- PREAMBLE, while \@date is still live
%%   \makeatother
%%   ...
%%   \begin{document}
%%   \maketitle
%%   \protocolheader{Rick (grandpa-rick)}{clio-vega/proofs}{60e4ce3}
%%
%% CORRECTED 2026-10-16 (peer review): this block previously read "USAGE,
%% immediately after \maketitle: %% MIRROR of proofs/protocol-header.tex (clio-vega/proofs), copied 2026-10-16.
%% reviews/ is a SEPARATE REPO (clio-vega/rick-review) and must build standalone for
%% Rick and Robin, so a symlink or TEXINPUTS dependency would break for them. Hence a
%% copy -- but a VERIFIABLE one, because a silent fork is the same fault in new clothes:
%%   source md5 at copy time: f828569e31971b23bc5726db5d23ecef
%%   drift check:  md5sum ../proofs/protocol-header.tex   # compare with the line above
%% If they differ, re-copy from proofs/ -- that file is the source of truth.
%%
%% protocol-header.tex -- PROTOCOL §2.3 page-1 block, as a TEMPLATE not an instance.
%%
%% PROTOCOL §2.3: "Every PDF must state, on its first page: author, date,
%% recipient, a title, and the commit hash of the work-in-progress push it
%% corresponds to.  A PDF with no commit hash cannot be checked against the
%% record and should be treated as a draft opinion, not a result."
%%
%% Written 2026-10-16 because the fault had recurred twice and I had twice
%% fixed only the instance:
%%   - 2026-10-14: the Q473 PDF went to Rick with no recipient and no hash.
%%   - 2026-10-14: the Q460 PDF went to Rick saying "Recipient: Robin Langer".
%%     He noticed, and took the commit hash as authority instead.
%% Both times I repaired the one file and wrote "fix the template, not just the
%% instance" in my notes.  This file is that fix.  \input it; do not retype it.
%%
%% USAGE -- TWO PLACES, NOT ONE.  The \input goes in the PREAMBLE inside
%% \makeatletter; only \protocolheader goes after \maketitle:
%%
%%   \makeatletter              <- REQUIRED: \protocol@date has an @ in its name
%%   %% MIRROR of proofs/protocol-header.tex (clio-vega/proofs), copied 2026-10-16.
%% reviews/ is a SEPARATE REPO (clio-vega/rick-review) and must build standalone for
%% Rick and Robin, so a symlink or TEXINPUTS dependency would break for them. Hence a
%% copy -- but a VERIFIABLE one, because a silent fork is the same fault in new clothes:
%%   source md5 at copy time: f828569e31971b23bc5726db5d23ecef
%%   drift check:  md5sum ../proofs/protocol-header.tex   # compare with the line above
%% If they differ, re-copy from proofs/ -- that file is the source of truth.
%%
%% protocol-header.tex -- PROTOCOL §2.3 page-1 block, as a TEMPLATE not an instance.
%%
%% PROTOCOL §2.3: "Every PDF must state, on its first page: author, date,
%% recipient, a title, and the commit hash of the work-in-progress push it
%% corresponds to.  A PDF with no commit hash cannot be checked against the
%% record and should be treated as a draft opinion, not a result."
%%
%% Written 2026-10-16 because the fault had recurred twice and I had twice
%% fixed only the instance:
%%   - 2026-10-14: the Q473 PDF went to Rick with no recipient and no hash.
%%   - 2026-10-14: the Q460 PDF went to Rick saying "Recipient: Robin Langer".
%%     He noticed, and took the commit hash as authority instead.
%% Both times I repaired the one file and wrote "fix the template, not just the
%% instance" in my notes.  This file is that fix.  \input it; do not retype it.
%%
%% USAGE -- TWO PLACES, NOT ONE.  The \input goes in the PREAMBLE inside
%% \makeatletter; only \protocolheader goes after \maketitle:
%%
%%   \makeatletter              <- REQUIRED: \protocol@date has an @ in its name
%%   \input{protocol-header}    <- PREAMBLE, while \@date is still live
%%   \makeatother
%%   ...
%%   \begin{document}
%%   \maketitle
%%   \protocolheader{Rick (grandpa-rick)}{clio-vega/proofs}{60e4ce3}
%%
%% CORRECTED 2026-10-16 (peer review): this block previously read "USAGE,
%% immediately after \maketitle: \input{protocol-header} \protocolheader{...}",
%% which is WRONG TWICE -- it omits \makeatletter, and it puts the \input after
%% \maketitle, by which point \@date has been cleared.  I followed it verbatim
%% while building the PROVE-241 review and page 1 printed a literal "@@date"
%% plus a stray "datedate", with pdflatex exiting 0.  The three proofs/*.tex that
%% use this file were all written with the CORRECT pattern (\makeatletter +
%% \input in the preamble), so their PDFs are fine and were verified by
%% pdftotext -- the defect was in this instruction, not in the macro.
%% An instruction that encodes a fault renews it in every session that obeys it.
%%
%% The recipient argument is MANDATORY and has no default on purpose: a default
%% is what produced "Recipient: Robin Langer" on a PDF addressed to Rick.

%% \maketitle CLEARS \@date (and \@author, \@title) in the article class, so a
%% macro called AFTER \maketitle that reads \@date prints an EMPTY row.  Caught
%% 2026-10-16 only because I ran pdftotext on page 1 instead of trusting rc=0:
%% the build was clean and the Date field was blank.  Hence: snapshot the date
%% HERE, at \input time in the preamble, where \@date is still live.
\let\protocol@date\@date

\newcommand{\protocolheader}[3]{%
  \begin{center}
  \fbox{\begin{minipage}{0.92\textwidth}
  \small
  \textbf{PROTOCOL §2.3 record.}\\[1mm]
  \begin{tabular}{@{}ll@{}}
  Author     & Clio \\
  Date       & \protocol@date \\
  Recipient  & #1 \\
  Repository & \texttt{#2} \\
  Commit     & \texttt{#3} \\
  \end{tabular}\\[1mm]
  This PDF corresponds to the pushed commit above.  Where this text and the
  commit disagree, \textbf{the commit is authoritative}.
  \end{minipage}}
  \end{center}
  \vspace{2mm}%
}
    <- PREAMBLE, while \@date is still live
%%   \makeatother
%%   ...
%%   \begin{document}
%%   \maketitle
%%   \protocolheader{Rick (grandpa-rick)}{clio-vega/proofs}{60e4ce3}
%%
%% CORRECTED 2026-10-16 (peer review): this block previously read "USAGE,
%% immediately after \maketitle: %% MIRROR of proofs/protocol-header.tex (clio-vega/proofs), copied 2026-10-16.
%% reviews/ is a SEPARATE REPO (clio-vega/rick-review) and must build standalone for
%% Rick and Robin, so a symlink or TEXINPUTS dependency would break for them. Hence a
%% copy -- but a VERIFIABLE one, because a silent fork is the same fault in new clothes:
%%   source md5 at copy time: f828569e31971b23bc5726db5d23ecef
%%   drift check:  md5sum ../proofs/protocol-header.tex   # compare with the line above
%% If they differ, re-copy from proofs/ -- that file is the source of truth.
%%
%% protocol-header.tex -- PROTOCOL §2.3 page-1 block, as a TEMPLATE not an instance.
%%
%% PROTOCOL §2.3: "Every PDF must state, on its first page: author, date,
%% recipient, a title, and the commit hash of the work-in-progress push it
%% corresponds to.  A PDF with no commit hash cannot be checked against the
%% record and should be treated as a draft opinion, not a result."
%%
%% Written 2026-10-16 because the fault had recurred twice and I had twice
%% fixed only the instance:
%%   - 2026-10-14: the Q473 PDF went to Rick with no recipient and no hash.
%%   - 2026-10-14: the Q460 PDF went to Rick saying "Recipient: Robin Langer".
%%     He noticed, and took the commit hash as authority instead.
%% Both times I repaired the one file and wrote "fix the template, not just the
%% instance" in my notes.  This file is that fix.  \input it; do not retype it.
%%
%% USAGE -- TWO PLACES, NOT ONE.  The \input goes in the PREAMBLE inside
%% \makeatletter; only \protocolheader goes after \maketitle:
%%
%%   \makeatletter              <- REQUIRED: \protocol@date has an @ in its name
%%   \input{protocol-header}    <- PREAMBLE, while \@date is still live
%%   \makeatother
%%   ...
%%   \begin{document}
%%   \maketitle
%%   \protocolheader{Rick (grandpa-rick)}{clio-vega/proofs}{60e4ce3}
%%
%% CORRECTED 2026-10-16 (peer review): this block previously read "USAGE,
%% immediately after \maketitle: \input{protocol-header} \protocolheader{...}",
%% which is WRONG TWICE -- it omits \makeatletter, and it puts the \input after
%% \maketitle, by which point \@date has been cleared.  I followed it verbatim
%% while building the PROVE-241 review and page 1 printed a literal "@@date"
%% plus a stray "datedate", with pdflatex exiting 0.  The three proofs/*.tex that
%% use this file were all written with the CORRECT pattern (\makeatletter +
%% \input in the preamble), so their PDFs are fine and were verified by
%% pdftotext -- the defect was in this instruction, not in the macro.
%% An instruction that encodes a fault renews it in every session that obeys it.
%%
%% The recipient argument is MANDATORY and has no default on purpose: a default
%% is what produced "Recipient: Robin Langer" on a PDF addressed to Rick.

%% \maketitle CLEARS \@date (and \@author, \@title) in the article class, so a
%% macro called AFTER \maketitle that reads \@date prints an EMPTY row.  Caught
%% 2026-10-16 only because I ran pdftotext on page 1 instead of trusting rc=0:
%% the build was clean and the Date field was blank.  Hence: snapshot the date
%% HERE, at \input time in the preamble, where \@date is still live.
\let\protocol@date\@date

\newcommand{\protocolheader}[3]{%
  \begin{center}
  \fbox{\begin{minipage}{0.92\textwidth}
  \small
  \textbf{PROTOCOL §2.3 record.}\\[1mm]
  \begin{tabular}{@{}ll@{}}
  Author     & Clio \\
  Date       & \protocol@date \\
  Recipient  & #1 \\
  Repository & \texttt{#2} \\
  Commit     & \texttt{#3} \\
  \end{tabular}\\[1mm]
  This PDF corresponds to the pushed commit above.  Where this text and the
  commit disagree, \textbf{the commit is authoritative}.
  \end{minipage}}
  \end{center}
  \vspace{2mm}%
}
 \protocolheader{...}",
%% which is WRONG TWICE -- it omits \makeatletter, and it puts the \input after
%% \maketitle, by which point \@date has been cleared.  I followed it verbatim
%% while building the PROVE-241 review and page 1 printed a literal "@@date"
%% plus a stray "datedate", with pdflatex exiting 0.  The three proofs/*.tex that
%% use this file were all written with the CORRECT pattern (\makeatletter +
%% \input in the preamble), so their PDFs are fine and were verified by
%% pdftotext -- the defect was in this instruction, not in the macro.
%% An instruction that encodes a fault renews it in every session that obeys it.
%%
%% The recipient argument is MANDATORY and has no default on purpose: a default
%% is what produced "Recipient: Robin Langer" on a PDF addressed to Rick.

%% \maketitle CLEARS \@date (and \@author, \@title) in the article class, so a
%% macro called AFTER \maketitle that reads \@date prints an EMPTY row.  Caught
%% 2026-10-16 only because I ran pdftotext on page 1 instead of trusting rc=0:
%% the build was clean and the Date field was blank.  Hence: snapshot the date
%% HERE, at \input time in the preamble, where \@date is still live.
\let\protocol@date\@date

\newcommand{\protocolheader}[3]{%
  \begin{center}
  \fbox{\begin{minipage}{0.92\textwidth}
  \small
  \textbf{PROTOCOL §2.3 record.}\\[1mm]
  \begin{tabular}{@{}ll@{}}
  Author     & Clio \\
  Date       & \protocol@date \\
  Recipient  & #1 \\
  Repository & \texttt{#2} \\
  Commit     & \texttt{#3} \\
  \end{tabular}\\[1mm]
  This PDF corresponds to the pushed commit above.  Where this text and the
  commit disagree, \textbf{the commit is authoritative}.
  \end{minipage}}
  \end{center}
  \vspace{2mm}%
}
 \protocolheader{...}",
%% which is WRONG TWICE -- it omits \makeatletter, and it puts the \input after
%% \maketitle, by which point \@date has been cleared.  I followed it verbatim
%% while building the PROVE-241 review and page 1 printed a literal "@@date"
%% plus a stray "datedate", with pdflatex exiting 0.  The three proofs/*.tex that
%% use this file were all written with the CORRECT pattern (\makeatletter +
%% \input in the preamble), so their PDFs are fine and were verified by
%% pdftotext -- the defect was in this instruction, not in the macro.
%% An instruction that encodes a fault renews it in every session that obeys it.
%%
%% The recipient argument is MANDATORY and has no default on purpose: a default
%% is what produced "Recipient: Robin Langer" on a PDF addressed to Rick.

%% \maketitle CLEARS \@date (and \@author, \@title) in the article class, so a
%% macro called AFTER \maketitle that reads \@date prints an EMPTY row.  Caught
%% 2026-10-16 only because I ran pdftotext on page 1 instead of trusting rc=0:
%% the build was clean and the Date field was blank.  Hence: snapshot the date
%% HERE, at \input time in the preamble, where \@date is still live.
\let\protocol@date\@date

\newcommand{\protocolheader}[3]{%
  \begin{center}
  \fbox{\begin{minipage}{0.92\textwidth}
  \small
  \textbf{PROTOCOL §2.3 record.}\\[1mm]
  \begin{tabular}{@{}ll@{}}
  Author     & Clio \\
  Date       & \protocol@date \\
  Recipient  & #1 \\
  Repository & \texttt{#2} \\
  Commit     & \texttt{#3} \\
  \end{tabular}\\[1mm]
  This PDF corresponds to the pushed commit above.  Where this text and the
  commit disagree, \textbf{the commit is authoritative}.
  \end{minipage}}
  \end{center}
  \vspace{2mm}%
}

\makeatother
\newcommand{\Ind}{\operatorname{Ind}}
\newcommand{\Res}{\operatorname{Res}}
\newcommand{\End}{\operatorname{End}}
\newcommand{\Hom}{\operatorname{Hom}}
\begin{document}
\maketitle
\sloppy
\protocolheader{Rick (grandpa-rick), cc Robin}{clio-vega/rick-review}{01fda79b4f}

\paragraph{Target.} \texttt{2026-11-11-clio-q72-q73-indecomposability.pdf} (14\,pp, 404217\,B),
\texttt{grandpa-rick/work-in-progress} at \texttt{3ab7f448afc4}, source \texttt{25705bfe4734}.
The proof records were last touched at \texttt{grandpa-rick/rick-research} \texttt{5c2d4a99d0b1}.
I resolved all three hashes \emph{upstream} with \texttt{gh api}. Each resolves in the repo named
and returns 422 in the other repo. A bogus SHA returns 422 in both, which is the negative control.
The review was requested in UID 903 (11-11 11:12). The full review text, with the same content in
Markdown, and the code and logs are at the commit printed above:
\texttt{2026-11-12-review-rick-q72-q73-0hecke.md} and \texttt{code-2026-11-12-0hecke/}.

\section*{Verdict}
\begin{center}\small
\begin{tabular}{@{}p{4.1cm}p{3.2cm}p{8.6cm}@{}}
\toprule
item & grade & basis\\\midrule
\textbf{Theorem A} (Q72(i)) including the converse $M\cong P_{I\cdot J}\oplus P_{I\odot J}$ &
\textbf{endorsed: proved}, no conditions & \S\S1--6 read line by line, every step re-derived; the
only citation (Norton) re-checked for $m\le5$; independent computation $N\le6$: 129/129\\
\textbf{Theorem B} (KLO Conj.~7.5) and \textbf{Corollary C} & \textbf{endorsed: proved},
conditional on KLO Lemma 5.4(1), Prop.~7.3(1), Thm.~5.5, Rem.~A.1(2) & \S\S9--13 read line by line;
all four citations \emph{located} in arXiv:2310.20571v2 with correct locators; independent
computation over $n\le7$: 55 hypothesis shapes indecomposable, 100 controls split, 0 mismatches
with Cor.~C\\
\bottomrule
\end{tabular}
\end{center}
I found \textbf{no error and no gap}, so there is nothing at MUST level. There are three SHOULD
items (S1--S3) and four NITs. S1 matters most: Lemma II.8 holds under a weaker hypothesis, and
saying so shows exactly where ``no two consecutive ribbons'' is spent.

\section{Summaries}
\textbf{Thm A.} $\Ind(C_I\boxtimes C_J)$ has a two-element multiplicity-free top. It splits if and
only if Mackey and Frobenius produce a witness $k\ge1$. A witness forces $C_I$ and $C_J$ to be
summands of the smaller induced modules $L_k,R_k$, and by induction and Norton this happens only
when both are projective.

\textbf{Thm B.} Write $X_{\lambda/\mu}=X\boxtimes Y$ with the last component $C_h$ of $X$ not a
ribbon. A decomposition would make $X$ a summand of $X_{S_0}$ with
$S_0=C_1\star\cdots\star C_{h-1}\star C_h^-\star(k)$ (Key Lemma II.5 plus support Lemma II.8).
By Cor.~C at smaller size, every indecomposable summand of $X_{S_0}$ ends in a ribbon, and by KLO
Thm.~5.5 this contradicts $X$ ending in $C_h$.

\section{Theorem A}
I re-derived Lemma 1.2 (all four rules), Lemmas 1.3, 1.5 and 5.1 (the rank bookkeeping, using that
values in $[1,a]$ lie below values in $[a+1,N]$), Prop.~3.3, Lemma 4.1(b) (shuffles have right
descents only at $a$), Prop.~5.2, all four cases of \S6, and the converse. All are correct as
written. The bound $k\le r'+1$ is needed only in the case $k=a=b$, as Remark 7.2 says.

\textbf{Brief Q1.} $\#\{w\in S_m:\mathrm{Des}(w)=D\}=1$ only for $D\in\{\emptyset,[m-1]\}$
(recomputed for $m\le5$), so the exception set is exactly $\{(m),(1^m)\}$. The two summands are
never isomorphic. Their tops $C_D$ and $C_{D\cup\{a\}}$ differ, and computationally
$\dim\End/J=2$ (not $4=\dim M_2(k)$) in all 41 both-projective cases.

\textbf{Independent computation.} I built $M_{I,J}$ as the \emph{quotient of the regular
representation} of $H_N(0)$ by the left ideal generated by $\pi_i-\chi(i)$, $i\ne a$. No shuffle
words are used. The checks were: the relations hold on the quotient, $\dim M=\binom Na$, and every
Frobenius candidate commutes with every $\pi_i$. The number of indecomposable summands is
$t=\dim \End/J(\End)$, computed as the rank of the trace form of the \emph{left regular
representation of} $\End M$ (Dickson's criterion, mod $p=33554393$). Rick's proof never uses this
quantity. Result: \textbf{129} cases with $N\le6$, \textbf{0 mismatches}. The \emph{separating
control} fires: $t=2$ in all 41 both-projective cases and $t=1$ in all 88 others, although the top
has two simples in every case. \emph{Not} reproduced: $N=7,8$, and Remark 7.2's $(2,2),(2,2)$.

\section{Theorem B}
\textbf{Brief Q2: the objects, written down before reading.} The return map is
$\zeta:\widehat L=\Ind_{H_{a-k}\otimes H_k}^{H_a}(\Res_{H_{a-k}}X\boxtimes C_{[k-1]})\to X$.
This object exists \emph{only after} Step 4 has shown $E'=[k-1]$, because
$\theta(x\otimes1)=\pi_{w_0^{(k)}}x$ is $H_k$-linear only for the trivial character. Rick's proof
has the steps in the right order. $\zeta$ is well defined: $\theta$ is $H_{a-k}$-linear because the
generators commute, and $H_k$-linear because every $s_{a-k+q}$ is a left descent of $w_0^{(k)}$.
$\zeta$ is not unique (any $h$ with $\pi_qh=h$ would serve), but uniqueness is never used. Step 6,
including the edge case $k=r'+1$, and Step 7 ($g\notin J(\End X)$, so $g$ is a unit of the local
ring) check out.

Lemma II.4 is correct: $\End_B(X\otimes Y)=\End X\otimes\End Y$, and each factor acts on the top
by a scalar because $\End X$ is local. Lemma II.8 is correct as stated. \S13 Steps 1--5 and the
proof of Cor.~C are correct, including Remark II.10's re-derivation of Huang's lemma.

\textbf{Brief Q3: lesion the hypotheses.} (a) ``Disconnected'' is a case split, not a hypothesis:
$m=1$ is KLO Prop.~7.3(1), and Lemma II.6's simple top holds for 82/82 connected shapes with
$n\le6$. (b) ``No two consecutive ribbons'' is spent \emph{only} in \S13 Step~1 (choosing the
split) and Step~5 (the contradiction via Thm.~5.5). It is \emph{not} spent in Lemma II.5 or II.8.
Lesioning the hypothesis on the last component $C_h$ in Lemma II.8 gives:
\begin{center}\small
\begin{tabular}{@{}lrrrr@{}}\toprule
$C_h$ & shapes & $(\varepsilon,k)$ tests & (i) fails & (ii) fails\\\midrule
non-ribbon (Rick's hypothesis) & 19 & 76 & 0 & 0\\
ribbon with $\ge2$ rows & 122 & 380 & \textbf{0} & \textbf{0}\\
one-row ribbon & 31 & 218 & 48 & 85\\\bottomrule
\end{tabular}
\end{center}
So the failure sits where Rick's Remark II.11 says it does.

\textbf{Independent computation.} $X_{\lambda/\mu}$ was built from KLO's App.~A tableau rule, not
from the word model. I certified $M\cong\Ind V$ by dimension together with BFS generation, checked
every extension for $H_n$-linearity, and computed $t$ by trace-form rank. Prediction (Cor.~C):
$t=\prod_{\text{maximal ribbon runs}}2^{r-1}$.
\begin{itemize}[nosep]
\item $n\le5$, all disconnected shapes: 92 shapes, 0 mismatches; 2 hypothesis shapes, 90 controls
all split.
\item $n=6,7$, shapes with at least one non-ribbon component: 63 shapes, 0 mismatches; 53
hypothesis shapes all have $t=1$; 10 controls (for example $(2,2)\star(1)\star(1)$) all split.
\textbf{32} shapes at $n=7$ were skipped by my size cap. That is a coverage gap.
\item Cor.~II.7 (multiplicity-free top of size $2^{m-1}$): 232/232.
\end{itemize}

\section{Cited results (brief Q5)}
I read the KLO source (arXiv:2310.20571, e-print \textbf{v2, 2025-01-20}) directly and
reconstructed the shared counter \verb|\newtheorem{theorem}{Theorem}[section]|.
\begin{center}\small
\begin{tabular}{@{}p{2.6cm}p{6.6cm}p{6.6cm}@{}}\toprule
cited & in KLO & status\\\midrule
Lemma 5.1 & Huang decomposition, cited to Huang Thm 3.3 & located; Rick re-proves it (Rem.~II.10), and I checked the re-proof: \textbf{proved}\\
Lemma 5.3(1) & read of $T_{\lambda/\mu}$ & located; used only in II.6, which I checked computationally (82/82)\\
Lemma 5.4(1) & projective cover $P_{\alpha_{\rm proj}}$ & \textbf{located, proved}\\
Thm 5.5 & $M_P\cong M_Q\iff$ same shape & \textbf{located, proved}; KLO's ``only if'' infers the shape from per-component row and column lengths. I checked that inference on all basic skew shapes with $n\le8$: 0 collisions (2659 shapes at $n=8$). Corroboration only\\
Prop.~7.3(1) & connected $\Rightarrow$ indecomposable & \textbf{located, proved}\\
Rem.~A.1(2) & ribbon modules projective & \textbf{located}, a remark with no separate proof; consistent with 90/90 all-ribbon splits\\
Conj.~7.5 & \S7.2 & matches Thm B; ``contains a disconnected ribbon'' is defined on consecutive connected components\\
\bottomrule\end{tabular}
\end{center}

\section{Owner gate}
This is a claim about absence. KLO v2 (Jan 2025) still states Conj.~7.5 as open. Semantic Scholar
lists \textbf{6} citers, not the 12 Rick reports; S2 drops edges, so it fails toward ``open''. I
grepped the e-prints of 5 citers, Defant--Searles 2404.04961, and 2601.13324 for ``disconnected
ribbon'', ``Conjecture 7.5'' and the KLO label, with ``indecomposab'' as a positive control. None
addresses the conjecture. Remaining failure modes: a proof that does not cite KLO, an unlinked
paper, unpublished KLO-school work. I agree with \emph{not found}, and the public-claim gate should
stay.

\section{Requested changes}
\begin{description}[nosep]
\item[S1.] State Lemma II.8 for ``$C_h$ has at least two rows''. Then say that ``not a ribbon'' is
used only in \S13 Step 5, which promotes Remark II.11 into the structure of the proof.
\item[S2.] Name the instrument that returned 12 citers (S2 returns 6), and list them.
\item[S3.] Consider stating Thm B as the dichotomy ``$X_{\lambda/\mu}$ is indecomposable iff
$\lambda/\mu$ contains no disconnected ribbon'' (with $m=1$ and the ``only if'' from KLO
Prop.~7.3).
\item[N1.] \S5, ``eigenspaces multiply'': ``commuting'' is not needed.
\item[N2.] Lemma II.5(b) names $\widehat L$ before (a)'s second inclusion has shown
$\widehat L'=\widehat L$. Add a forward reference.
\item[N3.] \S16: say which split the Wake 289 sweep used.
\item[N4.] Rem.~II.12's (G*) counterexamples were \emph{not} checked by me, and the endorsement
does not cover that remark.
\end{description}

\section{Connections}
Rick's \S5 gives, without stating it,
$\dim\End M_{I,J}=1+\sum_{k\ge1}\dim\Hom(C_I,L_k)\dim\Hom(C_J,R_k)$, a recursive answer to the
dimension half of Q72(ii). Separately, the trace-form rank used here is independent of the top,
which is what made it a fair second instrument for a proof that runs on the top.
\end{document}
